%%% FIELD proof-four-coordinate-bias
Put \(L=\log(en)\) and \(\widetilde Y=Y/2\).  We invoke the established
construction
\(\mathtt{stat\_discrete\_ate\_heterogeneity\_frontier/
thm:robust\text{-}upper\text{-}construction\text{-}resolution}\) only for
its explicit clipped Chebyshev factorial estimator and its proved pilot and
approximation certificates.  Its hypotheses hold after this rescaling.  In
fact, the cell-mean envelope gives
\[
 \left|E_P[\widetilde Y\mid A=a,X=k]\right|
 =\frac{|\mu_{ak}|}{2}\le \frac M2.
\tag{1}
\]
Moreover, Cauchy--Schwarz and the conditional fourth-moment condition in the
definition of \(\mathcal P_{d,\epsilon,M}\) give
\[
 \begin{split}
 E_P\!\left[\left(\widetilde Y-\frac{\mu_{ak}}2\right)^2
       \middle|A=a,X=k\right]
 &\le
 \left\{E_P\!\left[\left(\widetilde Y-\frac{\mu_{ak}}2\right)^4
       \middle|A=a,X=k\right]\right\}^{1/2}\\
 &=\frac14
 \left\{E_P[(Y-\mu_{ak})^4\mid A=a,X=k]\right\}^{1/2}
 \le\frac{M^2}{4}.
 \end{split}
\tag{2}
\]
The fixed-overlap hypothesis is unchanged by rescaling, and independence of
the pilot and estimation folds follows from iid sampling.  Thus the cited
construction may be applied to \(\widetilde Y\), after which its output is
multiplied by two.

We next verify, rather than assume, the asserted four-coordinate form.  For a
cell \(k\), write
\[
 x_0=p_k(1-\pi_k),\qquad x_1=p_k\pi_k,\qquad
 z_a=x_a\mu_{ak},\qquad p_k=x_0+x_1.
\tag{3}
\]
The cited construction uses constants
\[
 K=\lfloor\alpha_0L\rfloor,\qquad
 \alpha_0=\min\left\{1,(64\log6)^{-1},(512D_6)^{-1}\right\},
 \qquad D_6=8\log(27/4),
\tag{4}
\]
the light-mass scale \(B=4096L/m_1\), where \(m_1\asymp n\) is the
estimation-fold size, and pilot threshold \(256L\).  Select the numerical
coefficient in \(L_n\) so that \(c_L\ge\alpha_0\); then \(K\le L_n\).
For the Chebyshev polynomial \(T_K\), define the polynomials, with their
continuous values at zero,
\[
 H_K(t)=\frac{1-T_K(1-2t)}{2K^2},\qquad
 E_K(t)=\frac{H_K(t)}t,\qquad
 G_K(t)=\frac{1-E_K(t)}t.
\tag{5}
\]
The identities \(T_K(1)=1\) and \(T'_K(1)=K^2\) show that both divisions in
(5) are exact polynomial divisions.  Since
\(T_K(\cos\vartheta)=\cos(K\vartheta)\), for \(0\le t\le1\),
\[
 0\le t\{1-tG_K(t)\}=H_K(t)\le K^{-2}.
\tag{6}
\]
Also, \(G_K\) has degree at most \(K-2\).  Hence
\[
 q_a(x_0,x_1)=\frac{x_0+x_1}{B}G_K(x_a/B)
\tag{7}
\]
has total degree at most \(K-1\le L_n-1\), and therefore
\[
 (x_0,z_0,x_1,z_1)\longmapsto
 z_1q_1(x_0,x_1)-z_0q_0(x_0,x_1)
 \in\mathscr Q_{L_n}.
\tag{8}
\]
This is the required four-coordinate rewrite on every dyadic slice of
\(\mathcal K_{\epsilon,M}(B)\).

Let \(\mathcal A\) be the pilot-good event from the cited construction.  On
\(\mathcal A\), every pilot-selected light cell obeys \(p_k\le B/4\).
For such a cell, fixed overlap yields \(x_a\ge\epsilon p_k\), while
\(0\le x_a/B\le1\).  Equations (3), (6), and the cell-mean envelope imply
\[
 \begin{split}
 \left|p_k\mu_{ak}-\frac{p_kz_a}{B}G_K(x_a/B)\right|
 &=p_k|\mu_{ak}|\frac{H_K(x_a/B)}{x_a/B}\\
 &\le \frac{MB}{\epsilon K^2}.
 \end{split}
\tag{9}
\]
Condition on the pilot sample.  An ordered distinct-index falling-factorial
term with one outcome mark is unbiased for the corresponding monomial
\(z_ap_kx_a^j\): after expansion over distinct indices, independence makes
its expectation the product of the one-observation mass and mark moments.
Linearity of this identity, first conditionally on the pilot and then
unconditionally, and summation of (9) over the two arms and at most \(d\)
selected cells give
\[
 \left|E_{Q_P^{(n)}}\!\left[
 \widehat T_{\mathrm L}
 -\sum_{k\in\widehat{\mathcal L}}
 p_k(\mu_{1k}-\mu_{0k})\right]\mathbf 1_{\mathcal A}\right|
 \le \frac{2MdB}{\epsilon K^2}.
\tag{10}
\]
The cited pilot certificate, applied under fixed overlap and
\(d\le c_\epsilon nL\), states after decreasing \(c_\epsilon\) if necessary
that
\[
 Q_P^{(n)}(\mathcal A^c)
 \le2d\exp(-c_\epsilon256L)\le C_\epsilon n^{-6}.
\tag{11}
\]
Both the factorial output and its target contribution are clipped at the
source envelope, so (11) contributes at most \(C_\epsilon Mn^{-6}\) to
(10).  The same calculation is precisely the cited pilot-misclassification
bound.  Since \(B\asymp L/n\), \(K\asymp L\), and \(L\ge1\), (10)--(11)
yield, after enlarging the constant for the finitely many rows for which
\(K<2\),
\[
 \sup_{P\in\mathcal P_{d,\epsilon,M}}
 \left|E_{Q_P^{(n)}}\!\left[
 \widehat T_{\mathrm L}
 -\sum_{k\in\widehat{\mathcal L}}
 p_k(\mu_{1k}-\mu_{0k})\right]\right|
 \le C_\epsilon M\frac{d}{n\log(en)}.
\tag{12}
\]
Finally, the Chebyshev recurrence has \(O(L_n^2)\) coefficients per cell;
the fixed four-coordinate bookkeeping, falling-factorial recurrence, zero
fallbacks, and clipping therefore require \(O(dL_n^2)\), hence
\(O(dL_n^3)\), arithmetic operations.  Rescaling the single mark by one half
and multiplying the final output by two preserves every expectation identity,
zero fallback, clipping rule, and pilot certificate.  This proves all parts
of the lemma.

%%% FIELD partial-global-upper
The proved light-approximation lemma supplies
\[
 \left|E\widehat T_{\mathrm L}
 -E\sum_{k\in\widehat{\mathcal L}}
 p_k(\mu_{1k}-\mu_{0k})\right|
 \le C_\epsilon M\frac{d}{n\log(en)}.
\]
The frozen heavy--light handle also gives an exact center and variance
identity on the all-heavy pilot event.  If explicit light and heavy radii
satisfied their allocated conditional tail bounds and had expected sizes
\(C_{\epsilon,\alpha}M\{n^{-1/2}+d/[n\log(en)]\}\) and
\(C_{\epsilon,\alpha}Mn^{-1/2}\), respectively, the triangle inequality and
the allocation of the noncoverage probability would prove the target, and
intersection with \([-2M,2M]\) would not enlarge length.  The frozen core
does not define radii with those properties, so this implication cannot be
turned into a proof of the stated interval theorem.

%%% FIELD partial-finite-precision
Let the exact-real interval be \([L,U]\).  If outward interval arithmetic
returns \(\widetilde L\le L\le U\le\widetilde U\) with
\(L-\widetilde L\le\delta_{\mathrm{num}}\) and
\(\widetilde U-U\le\delta_{\mathrm{num}}\), then deterministically
\[
 [L,U]\subseteq[\widetilde L,\widetilde U],\qquad
 (\widetilde U-\widetilde L)-(U-L)\le2\delta_{\mathrm{num}}.
\]
Thus coverage is monotone under the certified enlargement.  The banked
coefficient bound and the declared working precision control the arithmetic
magnitudes of the Chebyshev and factorial recurrences.  What is unavailable
is a defined monotone statistical score with a positive root-slope modulus;
without it, score-enclosure error cannot be converted into the asserted
endpoint error or bisection count.

%%% FIELD partial-frontier
The established same-class testing converse and the deterministic interval
\([-2M,2M]\) give the closed bracket
\[
 c_{\epsilon,\alpha}r_{n,d,M}
 \le\mathsf W^*_{n,d,\epsilon,M,\alpha}\le4M.
\]
The upper interval is measurable, has uniform coverage one, and has length
\(4M\).  The matching upper inequality of order \(r_{n,d,M}\) and its
finite-precision attainment would follow from the two still-unproved upstream
targets, but neither follows from the converse or the range interval.

%%% FIELD partial-local-linearization
Let
\[
 \mathcal H_n=
 \left\{\min_{a\in\{0,1\},\,1\le k\le d_n}
 N_{ak}^{(0)}>256\log(en)\right\}.
\]
The established all-heavy pilot lemma gives
\(Q_{P_n}^{(n)}(\mathcal H_n^c)=o(1)\).  On \(\mathcal H_n\), the frozen
heavy--light handle imposes the exact identities
\[
 (\widehat\tau^{\mathrm{HL}},\widehat V^{\mathrm{HL}})
 =(\widehat\tau^{\mathrm{OS}},\widehat V^{\mathrm{OS}}).
\]
Consequently the established saturated one-step benchmark transfers the
heavy--light center expansion, the center comparison, variance-ratio
consistency, and the centered Gaussian limit: equality on events whose
probabilities tend to one preserves each convergence-in-probability and
weak-convergence assertion.  No corresponding identity is imposed on
\(\widehat R^{\mathrm{HL}}\).  Hence the claimed expansion of that radius to
the Gaussian critical radius remains unproved.

%%% FIELD partial-endpoint-equivalence
On the balanced triangular class, the all-heavy identity and the two
established local results imply
\[
 \widehat\tau^{\mathrm{HL}}-\widehat\tau^{\mathrm{OS}}
 =o_{Q_{P_n}^{(n)}}(n^{-1/2}),
 \qquad
 \sqrt{\widehat V^{\mathrm{HL}}}
 -\sqrt{\widehat V^{\mathrm{OS}}}=o_{Q_{P_n}^{(n)}}(1).
\]
The second relation uses the positive lower and finite upper bounds on
\(V_{\mathrm{eff}}(P_n)/M^2\), so division and the square-root map are
legitimate.  Were the missing radius relation
\[
 \widehat R^{\mathrm{HL}}
 -z_{1-\alpha/2}\sqrt{\widehat V^{\mathrm{HL}}/n}
 =o_{Q_{P_n}^{(n)}}(n^{-1/2})
\]
available, each unprojected endpoint difference would be
\(o_{Q_{P_n}^{(n)}}(n^{-1/2})\).  The maps
\(x\mapsto\max\{-2M,x\}\) and \(x\mapsto\min\{2M,x\}\) are
one-Lipschitz, so projection would preserve that order, including at the
boundary.  This proves the projection reduction but not the missing radius,
global honesty, or numerical assertions.

%%% FIELD partial-fixed-alphabet
If \(d_n=d\) is fixed, then \(d=o(\sqrt n)\), the balanced lower bound remains
\(p_k\ge c_0/d\), and the efficient variance remains in
\([v_0M^2,v_1M^2]\).  The established saturated one-step benchmark therefore
gives its efficient-influence-function expansion, ratio-consistent variance,
and Gaussian limit.  The all-heavy pilot failure probability is exponentially
small, so the frozen exact override transfers the same center and variance
conclusions to the heavy--light outputs.  Fixed alphabet does not define the
heavy--light radius or imply its Gaussian critical-value expansion, and hence
does not complete the stated Wald-endpoint reduction.
